% GENERATED FILE. Edit canonical lessons, then run npm run generate:books.
% canonical-source-hash: fnv1a64:d877ca2246401527
% canonical-lessons: RU-C99-glue, RU-C99-scissors, RU-C99-alarm-clock, RU-C99-camera, RU-C99-purse

\chapter{Glue, Scissors, Alarm clock, Camera, Purse}
\label{ch:ru-glue-scissors-alarm-clock-camera-purse}
\hlchaptermodality{99}

\begin{chapteropening}
\textbf{By the end of this chapter:} \emph{I can say glue, scissors, an alarm clock, a camera and a purse.}

Complete the last lesson of chapter 99: I can say glue, scissors, an alarm clock, a camera and a purse.
\end{chapteropening}

\section[\texorpdfstring{\ru{клей} (kley)}{клей (kley)}]{\ru{клей} (kley) — glue}

\label{lesson:RU-C99-glue}



\begin{quote}
Before the new one: say the Russian for a nail, then the Russian for a hammer.
\end{quote}

\subsection*{You'll want to know: \ru{клей}}
\textbf{\ru{клей}} (\emph{kley}) — \textquotedblleft{}glue\textquotedblright{}.

\begin{cousinweb}
\textbf{\ru{клей}} $\leftarrow$ Proto-Slavic \emph{*kŭlějĭ} — inherited, a cousin of Greek \emph{kólla}, \textquotedblleft{}glue\textquotedblright{} (Vasmer), the source of \emph{collage}.
\end{cousinweb}

\begin{grammarlens}[title={What you've built}]
One more word for this chapter.
\end{grammarlens}

\subsection*{Your turn}
\begin{itemize}
\raggedright
  \item \emph{Say it:} \emph{kley}
  \item \emph{Say it:} \emph{kley}, once more
  \item \emph{Say it:} say \emph{kley}
  \item \emph{Recall:} say the Russian for a nail, then the Russian for a hammer, then say \emph{kley} again
\end{itemize}

\begin{tcolorbox}[breakable,colback=teal!4,colframe=teal!35!black,title={Before you move on}]
What does \ru{клей} mean? (\textquotedblleft{}Glue\textquotedblright{}.) Say it once more.
\end{tcolorbox}
\section[\texorpdfstring{\ru{ножницы} (nózhnitsy)}{ножницы (nózhnitsy)}]{\ru{ножницы} (nózhnitsy) — scissors}

\label{lesson:RU-C99-scissors}



\begin{quote}
Before the new one: say the Russian for a hammer, then the Russian for glue.
\end{quote}

\subsection*{You'll want to know: \ru{ножницы}}
\textbf{\ru{ножницы}} (\emph{nózhnitsy}) — \textquotedblleft{}scissors\textquotedblright{}.

\begin{cousinweb}
\textbf{\ru{ножницы}} is built on \emph{\ru{нож}}, \textquotedblleft{}a knife\textquotedblright{}, with \emph{-\ru{ницы}}, a plural ending for paired tools: a pair of blades.
\end{cousinweb}

\begin{grammarlens}[title={What you've built}]
One more word for this chapter.
\end{grammarlens}

\subsection*{Your turn}
\begin{itemize}
\raggedright
  \item \emph{Say it:} \emph{nózhnitsy}
  \item \emph{Say it:} \emph{nózhnitsy}, once more
  \item \emph{Say it:} say \emph{nózhnitsy}
  \item \emph{Recall:} say the Russian for a hammer, then the Russian for glue, then say \emph{nózhnitsy} again
\end{itemize}

\begin{tcolorbox}[breakable,colback=teal!4,colframe=teal!35!black,title={Before you move on}]
What does \ru{ножницы} mean? (\textquotedblleft{}Scissors\textquotedblright{}.) Say it once more.
\end{tcolorbox}
\section[\texorpdfstring{\ru{будильник} (budílnik)}{будильник (budílnik)}]{\ru{будильник} (budílnik) — an alarm clock}

\label{lesson:RU-C99-alarm-clock}



\begin{quote}
Before the new one: say the Russian for glue, then the Russian for scissors.
\end{quote}

\subsection*{You'll want to know: \ru{будильник}}
\textbf{\ru{будильник}} (\emph{budílnik}) — \textquotedblleft{}an alarm clock\textquotedblright{}.

\begin{cousinweb}
\textbf{\ru{будильник}} is \emph{\ru{будить}}, \textquotedblleft{}to wake someone\textquotedblright{}, with \emph{-\ru{льник}}, a suffix for a tool. \emph{\ru{будить}} makes \emph{\ru{бдеть}}, \textquotedblleft{}to be awake\textquotedblright{}, causative, and Vasmer relates it to Sanskrit \emph{bodháyati}, \textquotedblleft{}wakes, teaches\textquotedblright{}, the root of \emph{Buddha}, \textquotedblleft{}the awakened\textquotedblright{}.
\end{cousinweb}

\begin{grammarlens}[title={What you've built}]
One more word for this chapter.
\end{grammarlens}

\subsection*{Your turn}
\begin{itemize}
\raggedright
  \item \emph{Say it:} \emph{budílnik}
  \item \emph{Say it:} \emph{budílnik}, once more
  \item \emph{Say it:} say \emph{budílnik}
  \item \emph{Recall:} say the Russian for glue, then the Russian for scissors, then say \emph{budílnik} again
\end{itemize}

\begin{tcolorbox}[breakable,colback=teal!4,colframe=teal!35!black,title={Before you move on}]
What does \ru{будильник} mean? (\textquotedblleft{}An alarm clock\textquotedblright{}.) Say it once more.
\end{tcolorbox}
\section[\texorpdfstring{\ru{фотоаппарат} (fotoapparát)}{фотоаппарат (fotoapparát)}]{\ru{фотоаппарат} (fotoapparát) — a camera}

\label{lesson:RU-C99-camera}



\begin{quote}
Before the new one: say the Russian for scissors, then the Russian for an alarm clock.
\end{quote}

\subsection*{You'll want to know: \ru{фотоаппарат}}
\textbf{\ru{фотоаппарат}} (\emph{fotoapparát}) — \textquotedblleft{}a camera\textquotedblright{}.

\begin{cousinweb}
\textbf{\ru{фотоаппарат}} is \emph{\ru{фото}-}, from Greek \emph{phōs}, \emph{phōtós}, \textquotedblleft{}light\textquotedblright{}, with \emph{\ru{аппарат}}, \textquotedblleft{}a device\textquotedblright{}, from Latin \emph{apparātus}, \textquotedblleft{}equipment\textquotedblright{}.
\end{cousinweb}

\begin{grammarlens}[title={What you've built}]
One more word for this chapter.
\end{grammarlens}

\subsection*{Your turn}
\begin{itemize}
\raggedright
  \item \emph{Say it:} \emph{fotoapparát}
  \item \emph{Say it:} \emph{fotoapparát}, once more
  \item \emph{Say it:} say \emph{fotoapparát}
  \item \emph{Recall:} say the Russian for scissors, then the Russian for an alarm clock, then say \emph{fotoapparát} again
\end{itemize}

\begin{tcolorbox}[breakable,colback=teal!4,colframe=teal!35!black,title={Before you move on}]
What does \ru{фотоаппарат} mean? (\textquotedblleft{}A camera\textquotedblright{}.) Say it once more.
\end{tcolorbox}
\section[\texorpdfstring{\ru{кошелёк} (koshelyók)}{кошелёк (koshelyók)}]{\ru{кошелёк} (koshelyók) — a purse, a wallet}

\label{lesson:RU-C99-purse}



\begin{quote}
Before the new one: say the Russian for an alarm clock, then the Russian for a camera.
\end{quote}

\subsection*{You'll want to know: \ru{кошелёк}}
\textbf{\ru{кошелёк}} (\emph{koshelyók}) — \textquotedblleft{}a purse, a wallet\textquotedblright{}.

\begin{cousinweb}
\textbf{\ru{кошелёк}} is \emph{\ru{кошель}}, \textquotedblleft{}a bag\textquotedblright{}, with the diminutive \emph{-\ru{ёк}}. \emph{\ru{кошель}} comes from \emph{\ru{кош}}, \textquotedblleft{}a basket\textquotedblright{} (Vasmer), which he relates to Latin \emph{quālum}, \textquotedblleft{}a wicker basket\textquotedblright{}.
\end{cousinweb}

\begin{grammarlens}[title={What you've built}]
One more word for this chapter.
\end{grammarlens}

\subsection*{Your turn}
\begin{itemize}
\raggedright
  \item \emph{Say it:} \emph{koshelyók}
  \item \emph{Say it:} \emph{koshelyók}, once more
  \item \emph{Say it:} say \emph{koshelyók}
  \item \emph{Recall:} say the Russian for an alarm clock, then the Russian for a camera, then say \emph{koshelyók} again
\end{itemize}

\begin{tcolorbox}[breakable,colback=teal!4,colframe=teal!35!black,title={Before you move on}]
What does \ru{кошелёк} mean? (\textquotedblleft{}A purse, a wallet\textquotedblright{}.) Say it once more.
\end{tcolorbox}
